% =====================================================================
% Table tab:controls — the full control grid, cohort v4.
%
% Source: results_controls/{patient,site}_seed{42,142,242}/oof/oof_*.csv,
%   360 scored units (3 seeds x 2 protocols x 3 targets x 2 conditions
%   x 5 folds), plus results_tumor/{patient,site}_seed{42,142,242} for the
%   tumour-versus-normal arm, a further 60 units (3 seeds x 2 protocols
%   x 1 target x 2 conditions x 5 folds). Pooled out-of-fold AUROC, slides
%   aggregated to the patient by mean probability before scoring, exactly
%   as tab_mutation_luad does. Recomputed by scripts/verify_controls.py,
%   which writes controls_recomputed.csv.
%
% NOTE(tumour arm): added after the sex control returned an uninformative
%   result. It exists to break a circularity in the earlier draft, which
%   used TP53 both as the paper's headline measurement and as the proof
%   that the pipeline producing it works. Tumour-versus-normal is
%   independent of every mutation label and is the one discrimination H&E
%   is certain to support, so it can carry that burden without borrowing
%   from the arm it is meant to validate. 520 slides: the 458 of the main
%   cohort plus the 62 solid-tissue-normal sections Methods excludes.
%
% NOTE(seeds): the control arm ran three seeds, not the five of the main
%   arm. The SD columns are therefore over three values and are reported
%   only to show spread, not as an interval. Inference against chance
%   uses the analytic Hanley-McNeil SE quoted in the caption.
%
% NOTE(supersedes): this table replaces the four loose numbers the earlier
%   draft carried in the Results text (0.515 / 0.482 / 0.588 / 0.647). Those
%   were not wrong, they were a different estimator: they are the mean of the
%   five fold AUROCs for seed 42 under the patient-preserved protocol, read
%   off results_controls/patient_seed42/summary_table.csv. tab_mutation_luad
%   reports pooled out-of-fold AUROC, so the controls are recomputed the same
%   way here. The two estimators differ by up to 0.05 on the same predictions
%   (sex: 0.588 as a mean of folds, 0.549 pooled, both for seed 42), which is
%   itself a reason to state which one a number is.
% =====================================================================
\small
\setlength{\tabcolsep}{5pt}
\renewcommand{\arraystretch}{1.15}
\begin{tabular}{l l r c c c c}
\toprule
\textbf{Target} & \textbf{Cond.} & $\boldsymbol{n_+}$ &
\textbf{s42} & \textbf{s142} & \textbf{s242} & \textbf{mean}$\pm$\textbf{SD} \\
\midrule
\multicolumn{7}{l}{\textbf{a}, Patient-preserved outer folds}\\
\addlinespace[1pt]
tumour vs.\ normal\textsuperscript{$\ddagger$}
                & B2 & $444$ & .980 & .972 & .977 & \textbf{.976}$\pm$.004 \\
tumour vs.\ normal\textsuperscript{$\ddagger$}
                & B3 & $444$ & .775 & .815 & .810 & \textbf{.800}$\pm$.022 \\
\addlinespace[2pt]
\textit{TP53}   & B2 & $220$ & .636 & .696 & .684 & .672$\pm$.032 \\
\textit{TP53}   & B3 & $220$ & .529 & .594 & .595 & .573$\pm$.037 \\
\addlinespace[2pt]
\textit{TP53} permuted & B2 & $220$ & .508 & .489 & .493 & .496$\pm$.010 \\
\textit{TP53} permuted & B3 & $220$ & .476 & .482 & .504 & .487$\pm$.015 \\
\addlinespace[2pt]
patient sex     & B2 & $238$ & .549 & .506 & .563 & .539$\pm$.029 \\
patient sex     & B3 & $238$ & .467 & .520 & .496 & .494$\pm$.027 \\
\midrule
\multicolumn{7}{l}{\textbf{b}, Site-preserved outer folds}\\
\addlinespace[1pt]
tumour vs.\ normal\textsuperscript{$\ddagger$}
                & B2 & $444$ & .956 & .982 & .965 & \textbf{.968}$\pm$.013 \\
tumour vs.\ normal\textsuperscript{$\ddagger$}
                & B3 & $444$ & .731 & .788 & .771 & \textbf{.763}$\pm$.029 \\
\addlinespace[2pt]
\textit{TP53}   & B2 & $220$ & .617 & .657 & .660 & .645$\pm$.024 \\
\textit{TP53}   & B3 & $220$ & .590 & .594 & .567 & .584$\pm$.014 \\
\addlinespace[2pt]
\textit{TP53} permuted & B2 & $220$ & .514 & .463 & .465 & .481$\pm$.029 \\
\textit{TP53} permuted & B3 & $220$ & .486 & .512 & .472 & .490$\pm$.021 \\
\addlinespace[2pt]
patient sex     & B2 & $238$ & .525 & .532 & .531 & .530$\pm$.004 \\
patient sex     & B3 & $238$ & .494 & .481 & .494 & .489$\pm$.008 \\
\bottomrule
\end{tabular}
\begin{flushleft}\footnotesize
Pooled out-of-fold \auroc{}, three seeds, both protocols. B2 and B3 are
the conditions of \cref{tab:mutation_luad_resultssec}. Labels are
permuted at the patient level, so a slide never carries a label drawn
from its own patient. \textsuperscript{$\ddagger$}The tumour-versus-normal
arm scores $506$ patients---the $444$ of the main cohort against the
$62$ solid-tissue-normal sections excluded from it (Methods)---and no
patient contributes to both classes, so the patient-level pooling is the
same operation as in every other row. All other rows score the $444$.
\par\medskip
\emph{The negative control passes.} All four permuted rows lie within
$0.019$ of chance, and every one contains $0.5$ in its analytic
$95\%$ interval. Nothing in the protocol manufactures discrimination.
\par\medskip
\emph{The positive control passes, on both channels.} Tumour versus
normal reaches $0.976$ patient-preserved (analytic $95\%$ CI
$0.965$--$0.988$, $81$ standard errors from chance) and $0.968$
site-preserved. This is the control that establishes the pipeline
learns, and it was chosen because it is the one discrimination H\&E is
certain to support, so a failure here would have been unambiguous. The
pattern channel clears chance too, at $0.800$ ($0.752$--$0.848$) and
$0.763$: the ANORAK head is functional, which forecloses the reading
that its near-chance mutation values in
\cref{tab:mutation_luad_resultssec} reflect a broken classifier rather
than an absent signal. Neither number licenses any inference about
mutation prediction. A cohort where the intended target is separable at
$0.98$ and \textit{TP53} stops at $0.66$ is a cohort in which the
machinery works and the target is hard, which is precisely the
distinction the controls exist to draw.
\par\medskip
\emph{\textit{TP53} replicates the main arm.} At $0.672$
patient-preserved it sits $0.011$ from the five-seed estimate of
\cref{tab:mutation_luad_resultssec} ($0.661\pm0.033$), and
site-preserved the gap is $0.005$ ($0.645$ vs.\ $0.640$). Three seeds
run for a different purpose therefore reproduce the main arm. We report
this as replication rather than as evidence that the pipeline learns,
which would use the measurement to certify the instrument that produced
it; the tumour arm carries that burden instead.
\par\medskip
\emph{The sex control is uninformative, and we say so.} Patient sex was
prespecified before the tumour arm existed, on the expectation that it
is recoverable from H\&E. It is not, here: $0.539$
($0.485$--$0.593$, $p\!=\!0.15$) on the embedding channel and $0.494$ on
the pattern channel, with the site-preserved arm lower still. The design
cannot settle this either way---at $n_+\!=\!238$ the smallest deviation
from chance it resolves at $80\%$ power is $0.077$, so any true sex
effect below $0.577$ was always going to return this result. It is
therefore neither passed nor failed but underpowered, which is the same
verdict this paper issues on the aggregator gains of
\cref{tab:literature_power_audit} and we see no principled reason to
exempt ourselves from it. One inference does survive: because sex is not
recoverable from these embeddings at this scale, the \textit{TP53} and
\textit{EGFR} signal in \cref{tab:mutation_luad_resultssec} cannot be a
demographic shortcut. A control that fails to establish what it was
chosen to establish can still exclude something else.
\end{flushleft}
